\chapter{Methodology}
\label{ch:method}

This chapter describes the methodology adopted to address the research objectives stated in Chapter~\ref{ch:into}. It details the multi-agent system architecture, environment configuration, custom reward shaping wrapper formulation, PPO shared-policy training pipeline, and the mathematical design of the advanced football analytics and spatial pitch heatmaps.

% ================================================================
\section{Multi-Agent System Architecture}
\label{sec:method_design}
% ================================================================

The methodology is designed to systematically solve cooperative-competitive continuous control on a 3D physical soccer pitch. The architectural flow is structured into four main phases:

\begin{enumerate}
    \item \textbf{Environment Wrapper Architecture:} Constructing unified parallel wrappers to integrate the physical MuJoCo simulation with Stable-Baselines3.
    \item \textbf{Custom Reward Shaping Wrapper:} Mitigating sparse rewards using continuous state-dependent shaped rewards.
    \item \textbf{Cooperative Shared-Policy Training:} Deploying actor-critic Proximal Policy Optimisation (PPO) with parameter sharing.
    \item \textbf{Soccer Analytics \& Heatmap Visualization:} Processing trajectory states to calculate possession, teammate spacing, and 2D spatial pitch coverage densities.
\end{enumerate}

% ================================================================
\section{MuJoCo Soccer 2v2 Simulation Environment}
\label{sec:method_soccer_env}
% ================================================================

The experimental environment is the DeepMind Control Suite's multi-agent soccer task~\citep{tassa2020dmcontrol}, accessed via the Shimmy compatibility layer~\citep{shimmy2022} and structured using the PettingZoo parallel API~\citep{terry2021pettingzoo}. The simulation implements a physically realistic 2-versus-2 football match within a continuous 3D arena.

\subsection{Physical and Kinematic Specifications}
The environment contains four embodied agents (two per team: Red/Home and Blue/Away). The physics dynamics are governed by the MuJoCo engine, which simulates realistic friction, rigid body collisions, joint limits, and momentum. The field dimensions are scaled within continuous coordinates representing approximately $[-12, 12]$ on the longitudinal X-axis (pitch length) and $[-8, 8]$ on the lateral Y-axis (pitch width).

\subsection{State and Action Space Formulations}
The environment features high-dimensional continuous spaces:
\begin{itemize}
    \item \textbf{State Observations}: Each agent receives a local egocentric observation dictionary:
    \begin{itemize}
        \item \texttt{ball\_ego\_position}: The relative XY coordinates of the ball in the agent's coordinate system.
        \item \texttt{stats\_vel\_to\_ball}: The relative scalar velocity of the agent moving toward the ball.
        \item \texttt{stats\_vel\_ball\_to\_goal}: The velocity of the ball relative to the opponent's goal mouth.
    \end{itemize}
    \item \textbf{Action space}: Continuous torques applied directly to the agent's motor joints, enabling fluid movements (running, turning, and kicking).
\end{itemize}

% ================================================================
\section{Proximal Policy Optimisation (PPO) with Shared Policies}
\label{sec:method_ppo}
% ================================================================

\subsection{Custom Reward Shaping}
\label{sec:rw_shaping}

The default environment provides only sparse goal-scoring rewards, which are 
too infrequent to drive learning in the early stages of training. A custom 
\texttt{SoccerRewardShapingWrapper} was implemented as a PettingZoo 
\texttt{BaseParallelWrapper}, augmenting the sparse reward with three 
shaped components at each time step:

\begin{equation}
    r_{\text{shaped}}(a) = r_{\text{env}}(a) + r_{\text{dist}}(a) 
    + r_{\text{vel}}(a) + r_{\text{goal}}(a),
    \label{eq:total_reward}
\end{equation}

where:
\begin{align}
    r_{\text{dist}}(a) &= -0.01 \times \|\mathbf{b}_{\text{ego}}\|_2 
    \label{eq:dist_penalty} \\
    r_{\text{vel}}(a)  &= 0.10 \times v_{\text{to\_ball}} 
    \label{eq:vel_reward} \\
    r_{\text{goal}}(a) &= 0.50 \times v_{\text{ball\_to\_goal}} 
    \label{eq:goal_reward}
\end{align}

Here, $\|\mathbf{b}_{\text{ego}}\|_2$ is the Euclidean distance from the 
agent to the ball in the XY plane; $v_{\text{to\_ball}}$ is the scalar velocity 
of the agent toward the ball; and $v_{\text{ball\_to\_goal}}$ is the scalar 
velocity of the ball toward the opponent's goal. 

The distance penalty (Eq.~\eqref{eq:dist_penalty}) discourages passive 
behaviour by penalising agents that stay far from the ball. The velocity 
reward (Eq.~\eqref{eq:vel_reward}) incentivises chasing the ball. The goal 
reward (Eq.~\eqref{eq:goal_reward}) provides a strong signal for productive 
ball advancement.

\subsection{Parameter Sharing and Environment Wrapping}

All four agents share a single PPO policy network. This is implemented by 
converting the PettingZoo parallel environment into a Stable-Baselines3 
\texttt{VecEnv} using SuperSuit~\citep{supersuit2021}:

\begin{lstlisting}[language=Python, 
    caption={Environment wrapping for parameter sharing.}, 
    label=list:env_wrapping]
env = create_soccer_env(render_mode=None)
env = ss.pettingzoo_env_to_vec_env_v1(env)
env = ss.concat_vec_envs_v1(env, num_vec_envs=1, 
                             num_cpus=1, 
                             base_class='stable_baselines3')
\end{lstlisting}

Each of the four agents' observations is treated as an independent sample, 
effectively creating four parallel training trajectories from a single 
environment rollout, all contributing gradient updates to one shared network.

\subsection{PPO Hyperparameters}

The PPO agent is instantiated via Stable-Baselines3~\citep{raffin2021stablebaselines3} with a \texttt{MultiInputPolicy} (required for dictionary observations). Hyperparameters are listed in Table~\ref{tab:ppo_hyperparams}.

\begin{table}[!ht]
    \centering
    \caption{PPO hyperparameters for MuJoCo Soccer 2v2.}
    \label{tab:ppo_hyperparams}
    \begin{tabular}{lc}
        \toprule
        \textbf{Hyperparameter} & \textbf{Value} \\
        \midrule
        Policy              & MultiInputPolicy \\
        Learning rate       & $3 \times 10^{-4}$ \\
        Batch size          & 4\,096 \\
        $n$\_steps          & 4\,096 \\
        Entropy coefficient & 0.01 \\
        Total timesteps     & 10\,000\,000 \\
        Framework           & Stable-Baselines3 \\
        \bottomrule
    \end{tabular}
\end{table}

\subsection{Training Infrastructure}

Training was conducted on Google Colaboratory using a T4 GPU (15 GB VRAM) 
provided free of charge. A \texttt{CheckpointCallback} was configured to 
save model checkpoints to Google Drive every 200\,000 steps, preventing 
data loss in the event of session termination. An anti-disconnection script 
was also used in the browser console to prevent the 90-minute inactivity 
timeout from interrupting training.

% ================================================================
\section{Ethical Considerations}
\label{sec:ethics}
% ================================================================

This project involves no human participants, no collection of personal data, 
and no sensitive information. All environments used are publicly available 
simulations. The following ethical dimensions are nonetheless considered:

\begin{itemize}
    \item \textbf{Research Integrity:} All code was written by the authors. 
    External libraries are explicitly cited. No plagiarism of code or text 
    was committed.
    
    \item \textbf{Environmental Impact:} GPU training on cloud infrastructure 
    carries an energy cost. Google Colab's shared GPU resources were used 
    responsibly and sessions were terminated when not required.
    
    \item \textbf{Reproducibility:} All hyperparameters, random seeds (where 
    applicable), and training configurations are documented in this report 
    and in the project codebase to ensure reproducibility.
    
    \item \textbf{Impact on Society:} Multi-agent reinforcement learning 
    has potential applications in robotics, autonomous systems, and 
    game AI. The methods developed here are educational in nature and 
    carry no direct negative societal impact.
\end{itemize}

% ================================================================
\section{Summary}
\label{sec:method_summary}
% ================================================================

This chapter described the full methodology of the project. The 2v2 continuous MuJoCo Soccer simulation environment, local egocentric observations, and custom reward shaping formulation were defined. The cooperative shared-policy PPO training pipeline using SuperSuit parameter sharing was detailed. Additionally, the training infrastructure on Google Colab and ethical considerations were documented. Chapter~\ref{ch:results} presents the qualitative and quantitative results obtained by applying this methodology.